\section{Related Work}

\paragraph{Diffusion-/Flow-based Models.}
Diffusion models (\eg, \citealt{sohl2015deep,ho2020denoising,song2020score}) and their flow-based counterparts (\eg, \citealt{lipman2022flow,liu2022flow,albergo2023stochastic}) formulate noise-to-data mappings through \textit{differential equations} (SDEs or ODEs). At the core of their inference-time computation is an \textit{\mbox{iterative}} update, \eg, of the form $\x_{i+1}=\x_{i} + \Delta\x_{i}$, such as with an Euler solver. The update $\Delta\x_{i}$ depends on the neural network $f$, and as a result, generation involves multiple steps of network evaluations.

A growing body of work has focused on reducing the steps of diffusion-/flow-based models.
Distillation-based methods 
(\eg, \citealt{salimans2022progressive,luo2023diff,yin2024one,zhou2024score}) distill a pretrained multi-step model into a single-step one.
Another line of research  aims to train one-step diffusion/flow models from scratch (\eg, \citealt{song2023consistency,frans2024one,boffi2025flow,geng2025mean}). To achieve this goal, these methods incorporate the SDE/ODE dynamics into training by approximating the induced trajectories.
In contrast, our work presents a conceptually different paradigm and does not rely on SDE/ODE formulations as in diffusion/flow models.

\paragraph{Generative Adversarial Networks (GANs).}
GANs \cite{goodfellow2014generative} are a classical family of models that train a generator by discriminating generated samples from real data.
Like GANs, our method involves a single-pass network $f$ that maps noise to data, whose ``goodness'' is evaluated by a loss function; however, unlike GANs, our method does not rely on adversarial optimization.

\paragraph{Variational Autoencoders (VAEs).}
VAEs~\cite{kingma2013auto} optimize the evidence lower bound (ELBO), which consists of a reconstruction loss and a KL divergence term.
Classical VAEs are one-step generators when using a Gaussian prior. Today's prevailing VAE applications often resort to priors learned from other methods, \eg, diffusion \cite{rombach2022high} or autoregressive models \cite{esser2021taming}, where VAEs effectively act as tokenizers.

\paragraph{Normalizing Flows (NFs).}
NFs~\cite{rezende2015variational,dinh2016density,zhai2024normalizing} learn mappings from data to noise and optimize the log-likelihood of samples.
These methods require invertible architectures and computable Jacobians. Conceptually, NFs operate as one-step generators at inference, with computation performed by the \mbox{inverse} of the network.

\paragraph{Moment Matching.} 
Moment-matching methods \cite{dziugaite2015training,li2015generative} seek to minimize the Maximum Mean Discrepancy (MMD) between the generated and data distributions.
Moment Matching has recently been extended to one-/few-step diffusion \cite{zhou2025inductive}.
Related to MMD, our approach also leverages the concepts of kernel functions and positive/negative samples. However, our approach focuses on a drifting field that explicitly governs the sample drifts at training time. Further discussion is in \ref{sec:theory_mmd_relation}.

\paragraph{Contrastive Learning.}
Our drifting field is driven by positive samples from the data distribution and negative samples from the generated distribution.
This is conceptually related to the positive and negative samples in \textit{contrastive representation learning} \cite{hadsell2006dimensionality,oord2018representation}.
The idea of contrastive learning has also been extended to generative models, \eg, to GANs~\cite{unterthiner2017coulomb,kang2020contragan} or Flow Matching~\cite{stoica2025contrastive}.